\section{Conclusion}
\label{sec:conclusion}

The economic consequences of advanced AI depend on more than whether machines can
perform existing jobs. A distinct transition occurs when AI can perform the research
required to produce better AI. The toy model represents this transition as a gradual
change in the composition of research. When human and automated-research services are
gross substitutes ($\sigma_{HM}>1$), the automated share rises with the research wage
relative to its unit resource cost.

Population growth matters while human research dominates. With $\phi<1$ and a
constant research-employment share, capability growth is semi-endogenous and equals
$\eta n/(1-\phi)$. Research automation creates a separate feedback from capability
to output, from output to automated research, and from research back to capability.
The fully Cobb--Douglas counterfactual isolates this channel. Its reduced dynamics
are stable when
$D_{\mathrm{CD}}=1-\phi-\eta\omega_M\omega_X/\omega_L>0$.
At $D_{\mathrm{CD}}=0$, capability is double exponential but remains finite at every
finite date. At $D_{\mathrm{CD}}<0$, the reduced system reaches a finite-time
singularity. Hence explosive growth does not require $\sigma_{XL}>1$.

When $\sigma_{HM}>1$, the automated contribution to research changes endogenously.
The relevant denominator becomes
$D(s^M_t)=1-\phi-\eta\upsilon s^M_t$ and declines as automated research gains weight.
Its AI-dominated limit,
$D_{\mathrm{AI}}=1-\phi-\eta\upsilon$, must be positive for the stated finite
constant-growth candidate. This condition is stricter than the human-essential one because the
entire growth of automated research, rather than only the Cobb--Douglas weight
$\omega_M$, feeds back into capability. A temporary negative value is not sufficient
for a singularity; the negative-denominator episode must be sufficiently persistent
to satisfy the integral condition derived in the paper.

These feedback results have a deliberately limited status. The baseline parameters
keep $D_{\mathrm{CD}}$ and its AI-dominated counterpart positive. The zero- and
negative-denominator cases are analytical counterfactuals under proportional
allocation. They do not establish that the fully optimizing monopoly follows the
same path: consumption, research, or another allocation can reach a boundary first.
Likewise, a failed constant-growth condition is not by itself proof of a finite-time
singularity.

The task-based extension identifies a separate extensive margin. Even when AI has a
lower marginal cost for a research task, a fixed integration and verification cost
can make human performance cheaper at low scale. Capability growth, a higher research
wage, or an expansion of the research program can move tasks across their automation
thresholds. The resulting aggregate elasticity can vary endogenously with the
distribution of task-level costs, while the CES remains the parsimonious baseline.

The production elasticity $\sigma_{XL}$ also separates the wage response from the
response of the production-labor share. An expansion of AI services reduces this share when
$\sigma_{XL}>1$, but it raises the wage when $\sigma_{XL}<1/\alpha$. Hence both outcomes occur
when $1<\sigma_{XL}<1/\alpha$. This static result is a benchmark, not the paper's
standalone contribution. In general equilibrium, capital deepening and the
reallocation of workers between production and research also affect wages and the
aggregate labor share. More precisely, aggregate labor-share growth equals
production labor-share growth plus the growth of population relative to production
labor. A rising human-research employment share raises the aggregate share relative
to the production share; a declining human-research share lowers it.

The capital-deepening benchmark clarifies which part of the growth result is specific
to recursive AI research. If the labor substitute is an ordinary reproducible capital
service, gross substitution can drive labor's share to zero while the wage rises, but
constant saving produces at most an asymptotically linear $AK$ technology with a
constant growth rate. In the baseline model, capability also lowers the unit resource
cost of AI services and is produced with AI research. This additional loop can make
the output--capital ratio and the return to capital increase without bound. The
common ingredient in the paper's acceleration mechanisms is therefore recursive
improvement, not gross substitution or capital deepening alone. A nonpositive
research-feedback denominator can generate acceleration at unit elasticities;
$\sigma_{XL}>1$ provides a distinct production-side amplifier when capability is
unbounded.

Research automation also has two distinct margins within the AI-development
industry. The ratio $H/M$ measures human intensity, while $H/N$ measures the scale
of human research employment. The first can fall while the second rises. This occurs
when automated research per capita expands more rapidly than human input per unit of
automated research contracts. In the equilibrium complementarity path, $H/N$
declines monotonically. In the Cobb--Douglas equilibrium, it rises during the first
part of the transition and later converges to zero while the automated research
share approaches one. The asymptotic disappearance of human research as a relative
input therefore does not imply that human research employment must fall throughout
the transition.

The production elasticity supplies the second dimension of the growth classification.
When $\sigma_{XL}<1$, the low-elasticity candidate leaves labor essential
and makes the productive return to further reductions in the unit cost of AI services
vanish. At $\sigma_{XL}=1$, positive wage growth would make automated research
asymptotically dominant. When $\sigma_{XL}>1$ and the effective unit cost of AI
services vanishes relative to capital intensity, the production-labor share converges
to zero. The unconditional nonexistence result is narrower and stronger where it
applies: no interior steady state can combine finite constant growth rates with a
positive constant capability-growth rate. The qualification is important:
$\sigma_{XL}>1$ alone does not exclude a boundary path with bounded capability.

The discontinuity at unit elasticity is an infinite-horizon classification, not a
jump in finite allocations. At every finite relative price, the CES contribution
share is continuous in $\sigma_{XL}$. But the limits $w/p^X\to\infty$ and
$\sigma_{XL}\downarrow1$ do not commute: the share remains $\omega_X$ at exactly
one and converges to one for every fixed elasticity above one. Under the paper's
high-capability regularity condition, unbounded capability induces this relative-price
limit. The relative-price threshold for the
transition diverges as $\sigma_{XL}$ approaches one. This makes the result a singular
perturbation and explains why an economy with $\sigma_{XL}$ only slightly above one
can look Cobb--Douglas over an extremely large range of finite relative prices
without possessing a finite-growth steady state asymptotically.

The numerical experiments solve the decentralized perfect-foresight canonical system.
Investment and research spending are endogenous, the household Euler equation holds,
and the developer satisfies its research first-order and costate conditions. With
$\sigma_{XL}=0.75$, output per capita falls initially, recovers, and then approaches
a constant level even though capability continues to grow. By year 600,
output-per-capita growth is 0.0056 percent and capability growth is 0.703 percent,
close to its analytical limit of 0.675 percent. The production-labor share rises as
labor becomes the bottleneck. With $\sigma_{XL}=1$ and $\sigma_{HM}=2$, capability
growth rises gradually and reaches 2.271 percent by year 2,000. Output- and
consumption-per-capita growth then approach 0.568 percent, while the automated
research share approaches one. Changing only the research elasticity to
$\sigma_{HM}=1$ lowers the limiting capability-growth rate to 1.738 percent and
per-capita output growth to 0.435 percent. The automated research contribution then
remains at 35 percent and human research employment converges to a positive share
of the population, even though $H/M$ converges to zero.
These transitions confirm that the asymptotic regime need not be a good medium-run
approximation.

For $\sigma_{XL}>1$, the paper solves a free-boundary equilibrium branch instead of
imposing a finite-growth terminal state. Conditional on an AI-dominated singular
limit, $g_A/(Y/K)$, the consumption share, the research-resource share, and $qA/K$
converge to constants, while $Y/K$ grows in inverse proportion to the remaining
time. The terminal conditions impose the limiting values of $C/Y$ and $qA/K$;
those two ratios are not independent validation moments. For $\sigma_{XL}=1.35$,
1.50, and 2.00, the non-imposed growth and resource-share ratios approach their
analytical predictions, while the initial jumps and estimated boundary dates remain
stable as the boundary moves outward. The estimated dates are 333.38, 280.45, and 209.08
years, respectively. These dates illustrate the comparative dynamics of the model;
they are not calibrated forecasts.

The return to capital supplies the price counterpart to these quantity dynamics.
Because saving is endogenous, the net return is an intertemporal equilibrium price
rather than a marginal product evaluated along an imposed investment path. It
converges to 4 percent under complementarity and to 4.568 percent in the
Cobb--Douglas AI-dominated limit. Under gross substitution, the net return first
remains moderate and then diverges with $Y/K$. Output and consumption growth also
accelerate. The real wage rises in all three reported gross-substitution paths,
while the aggregate labor share approaches zero. The earlier fixed-share paths
remain available as mechanism exercises but are not evidence about equilibrium
takeoff, interest rates, or wage declines.

A separate conditional search considers production elasticities above the static
wage threshold $1/\alpha$ and population growth as high as 3.9 percent. These paths
satisfy the canonical equilibrium equations on their computed domains but
extrapolate the AI-dominated terminal boundary beyond the region established by the
analytical result. Wage growth becomes negative in three of the four selected
configurations. The largest peak-to-trough wage decline is 9.39 percent when
$\sigma_{XL}=6.28$ and $n=3.9$ percent. The wage subsequently recovers and its trough
remains 16.72 percent above the initial wage. Capital deepening contributes to the
recovery, while the expansion of human employment in AI research makes production
labor scarce. More generally, the research first-order conditions imply that the
wage diverges whenever automated research per capita diverges. They also imply that,
with diminishing returns to the capability stock, a finite-time capability
singularity cannot occur while the wage remains bounded. The conditional singular
branch satisfies the stronger per-capita condition, so any finite wage decline on
that branch must eventually be reversed in levels. The high-elasticity wage-decline
searches use an extrapolated terminal boundary; proving recovery for every
infinite-horizon accelerating equilibrium still requires showing that acceleration
necessarily makes automated research per capita diverge.

Taken together, these results explain why the automation of AI research matters for
humans whether the real wage rises or falls temporarily. It can change the growth regime, reduce
labor's share of an expanding economy, increase the return to capital, and shift
income toward the owners of AI firms and other productive assets. It can also reduce
the number of humans required per unit of research without reducing human research
employment, because the development industry may expand faster than its human
intensity falls. The distinction between finite and asymptotic behavior is equally
important: observed stability over a finite range of capability and relative prices
does not by itself identify the long-run regime.

The current model does not establish whether an acceleration regime improves human
welfare. Monopoly rents accrue to a representative household, so the model abstracts
from unequal ownership and understates the distributional importance of a declining
labor share. It also excludes agency failures, misalignment, existential damage, and
political power. The numerical paths reported in the paper solve the decentralized
canonical equilibrium conditions. For $\sigma_{XL}>1$, they use the conditional singular
ratios derived in conditional result~\ref{cond:high-sigma-singular-scaling} and pass a
free-boundary convergence test. This procedure provides numerical evidence for a
locally regular singular branch from the stated initial conditions, but it is not a
proof of the developer's global dynamic optimality, global existence, or uniqueness.
The paper's acceleration and singularity calculations
remain conditional on the assumptions stated in the relevant conditional results and are
not quantitative forecasts. The paper does not assign probabilities to those
regimes or determine their full welfare ranking.

Market power creates a genuine policy trade-off. The monopolist restricts the current
supply of AI services, but its operating rents raise the private return to frontier
research. An explicit planner extension may value capability more than the monopolist
does; if it establishes $Q_t>q_t$, private research is insufficient at the private
allocation. The current model does not derive that inequality. Regulating the
AI-service price without replacing the lost innovation return can therefore improve
current deployment while slowing capability growth. The natural policy comparison is
between laissez-faire monopoly, regulated access, explicit research support, and a
joint policy that targets both margins.

A further extension should make explicit an agency problem that is only implicit in
the current model. As $H/M$ falls and automated systems choose a larger share of the
experiments, code, and model modifications that move the frontier, human owners and
researchers may become less able to observe their actions, verify their objectives,
or anticipate the consequences of every research decision. Supervision would then
be a costly choice rather than an implicit component of $H_t$. More intensive
oversight could reduce the probability that automated research departs from the
human principal's objectives, but it could also slow capability growth. Research
automation would thus generate a principal--agent problem with costly monitoring in
the spirit of \citet{holmstrom1979}.

Market structure may alter this trade-off, although its sign is not predetermined.
A protected monopolist internalizes a larger share of the loss to its installed
rents and franchise value if its automated researchers escape effective control. It
also faces less pressure to preempt a competitor. These forces may induce it to
moderate capability growth or invest more in supervision. Under monopolistic
competition, by contrast, each developer captures the private benefit of arriving
first, while part of the expected cost of inadequate supervision is borne by rivals
and society. Firms may consequently underinvest in oversight in order to accelerate
development and preserve their competitive position. Monopoly is not, however, a
guarantee of better control: concentration creates its own governance problems, and
the threat of entry or a foreign rival can restore the race incentive. The race and
slowdown branches in \citet{kokotajlo2025} illustrate precisely this tension: even a
concentrated frontier developer that observes evidence of misalignment may hesitate
to pause when another developer is close behind. A formal extension should therefore
allow developers to choose innovation speed and supervisory effort, make the
probability of an agency failure depend on automation and oversight, and compare
monopoly with monopolistic competition. If human research provides supervision that
is socially valuable but not fully remunerated, the extension could also evaluate a
second-best policy that subsidizes $H_t$ and taxes $M_t$. The purpose of these wedges
would not be to preserve human research employment, but to correct the effect of the
research-input mix on agency risk. A subsidy to verifiable supervision or liability
for agency failures would generally target the distortion more directly. A subsidy
to $H_t$ and a tax on $M_t$ become relevant when supervision cannot be separately
observed or contracted. These policies cannot be evaluated in the current model,
which contains neither agency losses nor an explicit supervision technology. The
central open question is when the
monopolist's internalization of agency losses dominates the governance costs of
concentration, and when competition turns effective supervision into an externality.
